% ======================================================================
% sections/conclusions.tex
%
% Conclusions for the full paper:
%   Triple Humanoid 24/7 Adverse Event Oncology Trial Response Team:
%   4-Site Rotation (v0.7.0)
%
% Four subsections inside the single \section{Conclusions} block.
% ======================================================================

\begin{sectionaccent}

\subsection{Thesis Restatement}
\label{sec:conclusions-thesis}

The thesis is preserved verbatim
\cite{kawchak_2026_humanoid_instructions}. On-premises repository
based LLMs provide commands to humanoid robots based on real-time
sensor data and controlled via x, y, z coordinates to administer
synergistic treatment regarding patient adverse events. This
workflow minimizes single robot error potential.

\subsection{Evidence Roll Up}
\label{sec:conclusions-evidence}

Table \ref{tbl:evidence-anchors} consolidates the 10 evidence anchors
that support the thesis. Each anchor is one row and
names a specific file in the read only RESEARCH trees so a reader
can audit the claim.

% \begin{figure}[H]
%     \hspace{1.75cm}
%    \includegraphics[width=0.8\linewidth, trim={0 0 0 0}, clip]{Images/07-08-process-funnel-chart.png}
%     % \vspace{-0.25cm}
%     \caption{New Caption}
%     \label{NewLabel}
% \end{figure}

\begin{table}[H]
\captionsetup{type=table}
\caption{Ten evidence anchors supporting the thesis. Each anchor
names a single project file. The list spans the swarm coordinator,
the FDA RTCT submitter, the 27 source file tree, the iteration
sweep index, the dashboard, the error checker and tests, the H2
dispatcher, the per site decision streams, and the DuckDB plus JSONL
data layer.}
\label{tbl:evidence-anchors}
\small
\begin{tabular}{>{\raggedright\arraybackslash}p{5cm}
                >{\raggedright\arraybackslash}p{10cm}}
\toprule
\textbf{Evidence Anchor} & \textbf{Function} \\
\midrule
swarm coordinator (191 lines) & Camaraderie logic that animates the 3 robot Unitree H2 EDU swarm at one site. The compact line count is itself evidence of clean coordination logic. \\
fda rtct submitter (67 lines) & Submits CTCAE grade 2 or higher records to the FDA RTCT pilot intake API within 1 hour per the HR 9505 SLA with ed25519 signing in production. \\
27 source files in codegen/src/ & Span Python, C++, Rust, YAML, and Markdown and together implement the full sensor to LLM to robot pipeline. \\
iterations index jsonl & 32 deterministic iteration sweeps across 4 trial sites, 5 axes. \\
dashboard html & Adds project interpretability for future paper readers without requiring raw JSON inspection. \\
check errors py & Provides LLM credibility in generating error checking code. \\
test xyz safety py & Provides LLM credibility in generating test code. \\
h2 dispatcher py & Per site H2 dispatch proficiency. \\
4 site llm decisions jsonl files & 1 Hz logs at scale for 4 sites (SF-01, SD-01, BO-01, AT-01). \\
aggregates duckdb + jsonl & DuckDB and JSONL proficiency in the same execution layer. \\
\bottomrule
\end{tabular}
\end{table}

\subsection{Industry Checklist Framing}
\label{sec:conclusions-checklist}

The paper presents the body of work as a checklist of items already
completed so that the next sponsor, the next site, the next IRB,
and the next regulator do not have to redo the groundwork. Table
\ref{tbl:industry-checklist} pulls the Inventory at a Glance table
from the instructions tree into a paper ready format. Each row is
marked Done so the next team has a clean hand off
\cite{kawchak_2026_humanoid_instructions}.

\begin{figure}[H]
    % \hspace{-0.75cm}
   \includegraphics[width=1\linewidth, trim={0 135 0 0}, clip]{Images/07-04-value-proposition-canvas.png}
    % \vspace{-0.25cm}
    \caption{Blue = Products/Customer Jobs, Red = Pain Relievers/Pains, Green = Gain Creators/Gains, Orange = Fit/Transit Reserve}
    \label{NewLabel}
\end{figure}

\begin{table}[H]
\captionsetup{type=table}
\caption{Industry checklist. The contents are drawn from the
Inventory at a Glance table in
\fp{paper/instructions/README.md}. Every item is marked Done so the
next team can build on top of it rather than rebuild it.}
\label{tbl:industry-checklist}
\small
\begin{tabular}{>{\raggedright\arraybackslash}p{5.2cm}
                >{\raggedright\arraybackslash}p{8cm}
                >{\raggedright\arraybackslash}p{2cm}}
\toprule
\textbf{Property} & \textbf{Value} & \textbf{Status} \\
\midrule
Network & PAT-NET-001 with 4 sites & Done \\
Robots per site & 3 Unitree H2 EDU at 39 DOF & Done \\
Robots total & 12 H2 across 4 sites & Done \\
LLM per site & 1 on-premises Claude Opus 4.7 1M & Done \\
LLM cadence & 1 Hz broadcast to all 3 robots & Done \\
Motion cadence & 10 Hz per active H2 & Done \\
Monitoring window & 168 hours (1 week) & Done \\
AE volume & 84 with 24 SAE per iteration & Done \\
Iterations & 32 deterministic sweeps & Done \\
E-stop latency & 5 ms swarm wide & Done \\
Per arm cumulative force & 15 N compressions, 5 N otherwise & Done \\
Inter robot minimum distance & 1.2 m at rest, 0.4 m at handoff & Done \\
Cumulative cross robot force & 22 N during 3 robot transfer & Done \\
Audit chain extension & SHA-256 per tick & Done \\
Camaraderie reduction & factor of 3 & Done \\
\bottomrule
\end{tabular}
\end{table}

\subsection{Forward Statement}
\label{sec:conclusions-forward}

The United States needs to remain Number 1 in the world regarding
patient safety, efficacy, and speed benefits in oncological adverse
event robot swarms for clinical trials. The paper closes with an
invitation to the FDA, the sponsors, the sites, and the IRBs to
advance the new field of Physical AI oncology trial surgical
procedures together. The prior December 2025 author paper on the
FDA Adverse Event Reporting System established the foundation for
LLM driven adverse event work \cite{kawchak_2025_18029100}. The
current paper extends that foundation into the coordinated 3 robot
24/7 adverse event response team described above. The next step
belongs to the agencies, the sponsors, the sites, and the IRBs, as depicted in \autoref{waterfall}.

\begin{figure}[H]
    % \hspace{-0.75cm}
   \includegraphics[width=1\linewidth, trim={0 0 0 0}, clip]{Images/07-05-financial-assessment-waterfall.png}
    % \vspace{-0.25cm}
    \caption{5-Year TCO waterfall regarding robot acquisition, service contracts, and per-site clusters.}
    \label{waterfall}
\end{figure}

\subsubsection*{Real life feasibility insight.} The pipeline in this
paper is small enough to audit, large enough to deploy at 4 sites,
and tight enough to honor the FDA RTCT 1 hour service level
agreement. The work is ready for the human governance step.

\sectiontable{Conclusions planning matrix}{%
Thesis & C1 & 1 & instructions & Verbatim restatement & 1 paragraph & No paraphrase & Done \\
Evidence & C2 & 2 & 10 anchors & Table roll up & 1 page & 10 rows & Done \\
Check. & C3 & 3 & Inventory & Done checklist & 1 page & Industry handoff & Done \\
Forward & C4 & 4 & 2025 FDA cite & US leadership & 1 paragraph & Final citation & Done \\
}

\end{sectionaccent}
